\chapter{Gán nhãn và phân bổ tần số}
\chapterlead{\emph{Radio labeling} (gán nhãn radio),
\emph{\(k\)-safe labeling} (gán nhãn an toàn \(k\)) và
\emph{frequency assignment} (phân bổ tần số) cùng chia sẻ một mô-đun: một lựa
chọn chính xác tạo ra \emph{forbidden interval} (khoảng cấm) trên miền của đối
tượng khác. Sự khác nhau nằm ở cách tạo khoảng cấm và cách định nghĩa hàm mục
tiêu.}

\index{radio labeling}\index{k-safe labeling}
\index{frequency assignment}\index{forbidden interval}
\index{arc-consistency preprocessing}

\flowdiagram{Giá trị chính xác}{Chuỗi thứ tự}{Khoảng cấm}
  {Độ trải hoặc số giá trị}{Nghiệm chứng và cận dưới}

\section{Hai văn liệu nguồn: gán nhãn đồ thị và phân bổ tần số}

Các bài toán trong chương đi vào SAT từ hai văn liệu khác nhau. Gán nhãn đồ
thị nghiên cứu các ánh xạ số trên đỉnh/cạnh dưới nhiều quy tắc khoảng cách;
gán nhãn radio là một nhánh trong hệ phân loại rất rộng đó
\parencite{gallian2024labeling}. Phân bổ tần số bắt nguồn từ mô hình nhiễu vô
tuyến: mỗi máy phát có miền, ràng buộc phân cách, đôi khi có giá trị gán trước,
phân cực, nhiều loại nhiễu và nhiều hàm mục tiêu.

Khảo cứu của Aardal và cộng sự phân loại FAP theo miền, mô hình nhiễu và hàm
mục tiêu, đồng thời tổng hợp MIP, cận dưới lý thuyết đồ thị, tối ưu ràng buộc
và phương pháp kinh nghiệm \parencite{aardal2007frequency}. Từ góc nhìn SAT,
điểm quan trọng là không đồng nhất ba hàm mục tiêu:
\begin{itemize}
  \item \textbf{độ trải nhỏ nhất}: giảm hiệu giữa tần số lớn nhất và nhỏ nhất;
  \item \textbf{số giá trị nhỏ nhất}: giảm số tần số phân biệt được dùng;
  \item \textbf{nhiễu/chi phí nhỏ nhất}: cho phép vi phạm hoặc mức nhiễu có
        trọng số.
\end{itemize}
Chúng có thể dùng cùng các ràng buộc khả thi nhưng cần biến hàm mục tiêu và lập
luận cận dưới khác nhau.

\begin{center}
\captionof{table}{Ba lớp bài toán gán nhãn và tần số.}
\label{tab:labeling-fap-classes}
\begin{tabularx}{\textwidth}{@{}>{\bfseries\sffamily}p{.2\textwidth}YYY@{}}
\toprule
Lớp & Miền & Ràng buộc & Hàm mục tiêu điển hình\\
\midrule
Gán nhãn radio & Liên tục \([1,B]\) & Từ khoảng cách đồ thị & Độ trải/nhãn lớn nhất nhỏ nhất\\
Đơn ánh an toàn \(k\) & Liên tục \([1,B]\) & Khác nhau + khoảng cách cạnh &
Nhãn lớn nhất nhỏ nhất\\
FAP & Rời rạc theo máy phát & Phân cách nhiễu &
Độ trải, số giá trị hoặc chi phí\\
\bottomrule
\end{tabularx}
\end{center}

Các công trình của nhóm trong chương được dùng để minh họa ba phép chuyển:
xung đột từng cặp sang biên khoảng; miền liên tục sang miền rời rạc; và cận độ
trải sang bộ đếm giá trị phân biệt. Ba phép chuyển này kết nối hệ phân loại FAP
và gán nhãn đồ thị với các mô-đun SAT có thể dùng lại.

\section{Bổ đề khoảng cấm}

Cho \(x_{u,a}=1\) khi đối tượng \(u\) nhận giá trị \(a\). Với đối tượng \(v\),
đặt
\[
  g_{v,t}=1\iff f(v)\ge t.
\]
Giả sử lựa chọn \(a\) cấm khoảng \([L(a),U(a)]\) cho \(v\). Điều kiện hợp lệ
là
\[
  f(v)<L(a)
  \quad\lor\quad
  f(v)>U(a).
\]
Trên miền nguyên liên tục, ràng buộc trở thành
\[
  \neg x_{u,a}
  \lor
  \neg g_{v,L(a)}
  \lor
  g_{v,U(a)+1}.
  \tag{11.1}\label{eq:boundary}
\]

\begin{lemma}[Mệnh đề biên]\label{lem:boundary}
Giả sử chuỗi thứ tự \(g_{v,t}\) đúng với ngữ nghĩa
\(f(v)\ge t\). Khi \(x_{u,a}=1\), mệnh đề
\eqref{eq:boundary} được thỏa khi và chỉ khi giá trị của \(v\) nằm ngoài
\([L(a),U(a)]\).
\end{lemma}

\begin{proof}
Nếu \(x_{u,a}=1\), literal đầu sai. Literal
\(\neg g_{v,L(a)}\) đúng khi \(f(v)<L(a)\); literal
\(g_{v,U(a)+1}\) đúng khi \(f(v)\ge U(a)+1\). Phép
tuyển vì thế tương đương với điều kiện nằm ngoài khoảng.
\end{proof}

Trên miền rời rạc \(D_v\), không thể dùng trực tiếp chỉ số \(L,U\). Ta thay
chúng bằng hai biên theo giá trị:
\[
  \lambda_v(a)=\min\set{q\mid D_{v,q}\ge L(a)},
\]
\[
  \rho_v(a)=\min\set{q\mid D_{v,q}>U(a)}.
\]
Các literal biên sau đó dùng \(\lambda_v(a)\) và \(\rho_v(a)\) trong chuỗi thứ
tự theo chỉ số miền.

\subsection{Quan hệ với Support Encoding}

Trong Direct Encoding/hỗ trợ, lựa chọn \(x_{u,a}\) yêu cầu \(v\) nhận một
trong các giá trị còn hỗ trợ:
\[
  \neg x_{u,a}
  \lor
  \bigvee_{b\in D_v:\ b\notin[L(a),U(a)]}x_{v,b}.
\]
Mệnh đề biên nén phép tuyển này thành hai ngưỡng khi tập hỗ trợ là hợp của một
tiền tố và một hậu tố. Đây không phải một mẹo riêng của radio/FAP; đó là phép
phân tích mệnh đề hỗ trợ dựa trên tính lồi của tập cấm. Nếu tập cấm có nhiều
khoảng rời nhau, cần nhiều cặp biên hoặc một
biểu diễn BDD/MDD khác.

\begin{figure}[tb]
\centering
\begin{tikzpicture}[satfig,x=11mm,y=8mm]
  \foreach \a in {1,...,9}{
    \ifnum\a<4
      \node[satbox,minimum width=9mm,fill=PaleBlue] (a\a) at (\a,0) {\(\a\)};
    \else
      \ifnum\a<7
        \node[satbox,minimum width=9mm,fill=PaleGold,draw=Gold] (a\a) at (\a,0)
          {\(\a\)};
      \else
        \node[satbox,minimum width=9mm,fill=PaleTeal] (a\a) at (\a,0) {\(\a\)};
      \fi
    \fi
  }
  \draw[Blue,very thick] (.55,-.75)--(3.45,-.75)
    node[midway,below,satlabel,text=Blue]{\(\neg g_{v,4}\): tiền tố hợp lệ};
  \draw[Gold,very thick] (3.55,.75)--(6.45,.75)
    node[midway,above,satlabel,text=Gold]{khoảng cấm \([4,6]\)};
  \draw[Teal,very thick] (6.55,-.75)--(9.45,-.75)
    node[midway,below,satlabel,text=Teal]{\(g_{v,7}\): hậu tố hợp lệ};
\end{tikzpicture}
\caption{Một khoảng cấm liên tục được nén thành hai literal biên:
\(\neg g_{v,4}\lor g_{v,7}\).}
\label{fig:forbidden-interval-boundaries}
\end{figure}

Phép so sánh lan truyền theo dõi đường suy luận của từng biểu diễn. Mệnh đề hỗ
trợ trực tiếp loại giá trị chính xác; mệnh đề biên suy qua chuỗi thứ tự và liên
kết giá trị chính xác--thứ tự. Hai phép biểu diễn đạt cùng mức cắt tỉa khi đặc tả của các
chuỗi đó cung cấp đủ các bước suy luận
\parencite{gent2002arc,bacchus2007gac}.

\begin{keyidea}{Mô-đun dùng lại}
\texttt{ExactValue} chọn \(a\); \texttt{OrderChain} biểu diễn miền của \(v\);
\texttt{ForbiddenInterval} chỉ cần một hàm trả về hai biên. Gán nhãn radio,
gán nhãn an toàn \(k\) và MO-FAP khác nhau ở hàm này, không khác nhau ở lôgic của
\cref{lem:boundary}.
\end{keyidea}

\section[Radio k-labeling]{Radio \(k\)-labeling: gán nhãn radio \(k\)}

\subsection{Mô hình}

Cho đồ thị liên thông \(G=(V,E)\), khoảng cách đường đi ngắn nhất \(d(u,v)\) và
tham số \(k\). Phép gán nhãn radio \(k\) là ánh xạ
\(f:V\to\mathbb{Z}_{>0}\) thỏa
\[
  \card{f(u)-f(v)}
  \ge
  \Delta_{uv},
  \qquad
  \Delta_{uv}=k+1-d(u,v),
\]
cho mọi cặp có \(\Delta_{uv}>0\). Mục tiêu là tối thiểu màu lớn nhất
\[
  \min\max_{v\in V}f(v).
\]
Trong bài toán gán nhãn radio chuẩn, người ta thường xét
\(k=\operatorname{diam}(G)\).

Khoảng cách đồ thị tạo một đồ thị ràng buộc dày hơn đồ thị gốc: mọi cặp có
\(d(u,v)\le k\) đều bị ràng buộc, với mức phân cách giảm khi khoảng cách đồ
thị tăng. Bước tìm đường đi ngắn nhất cho mọi cặp vì thế không chỉ dành cho bộ
bộ kiểm tra xác định toàn bộ cấu trúc liên kết của phép biểu diễn. Các cặp có
\(\Delta_{uv}\le0\) được bỏ, còn các đỉnh trong cùng quỹ đạo có thể hỗ trợ phá
đối xứng.

\begin{figure}[tb]
\centering
\begin{tikzpicture}[satfig,node distance=19mm]
  \node[satstate] (v1) {\(v_1\)};
  \node[satstate,right=of v1] (v2) {\(v_2\)};
  \node[satstate,right=of v2] (v3) {\(v_3\)};
  \node[satstate,right=of v3] (v4) {\(v_4\)};
  \draw[Ink,very thick] (v1)--(v2)--(v3)--(v4);
  \draw[Blue,dashed,bend left=35] (v1) to node[above,satlabel]{\(d=2\)} (v3);
  \draw[Blue,dashed,bend left=35] (v2) to node[above,satlabel]{\(d=2\)} (v4);
  \draw[Gold,dotted,very thick,bend right=48] (v1) to
    node[below,satlabel]{\(d=3\)} (v4);
  \node[below=4pt of v1,satlabel,text=Teal] {\(f=1\)};
  \node[below=4pt of v2,satlabel,text=Teal] {\(f=4\)};
  \node[below=4pt of v3,satlabel,text=Teal] {\(f=7\)};
  \node[below=4pt of v4,satlabel,text=Teal] {\(f=2\)};
\end{tikzpicture}
\caption{Một phép gán nhãn radio \(3\) khả thi cho \(P_4\). Hai đỉnh kề nhau
cần chênh ít nhất \(3\), cặp cách hai cạnh cần \(2\), và hai đầu mút cần \(1\).}
\label{fig:radio-p4}
\end{figure}

\subsection{SAT trực tiếp}

Với một cận trên \(B\), đặt
\[
  x_{v,\ell}=1\iff f(v)=\ell,
  \qquad \ell\in[1,B].
\]
Mỗi đỉnh nhận đúng một nhãn. Với cặp \(u,v\), mọi cặp nhãn
\((a,b)\) có \(\card{a-b}<\Delta_{uv}\) tạo xung đột
\[
  \neg x_{u,a}\lor\neg x_{v,b}.
\]
Direct Encoding có phép giải mã đơn giản nhưng số xung đột tăng theo độ rộng
khoảng cấm.

\subsection{SAT thứ tự}

Thêm
\[
  g_{v,\ell}=1\iff f(v)\ge\ell.
\]
Biến chính xác được xác định tại điểm chuyển:
\[
  x_{v,\ell}\iff
  g_{v,\ell}\land\neg g_{v,\ell+1}.
\]
Nếu \(u\) nhận \(a\), \(v\) không được nhận nhãn trong
\([a-\Delta_{uv}+1,a+\Delta_{uv}-1]\). Mệnh đề biên là
\[
  \neg x_{u,a}
  \lor
  \neg g_{v,a-\Delta_{uv}+1}
  \lor
  g_{v,a+\Delta_{uv}},
\]
sau khi bỏ các literal nằm ngoài miền. Số mệnh đề khoảng cách trở thành tuyến tính
theo \(B\) cho mỗi cặp đỉnh, thay vì phụ thuộc thêm vào \(\Delta_{uv}\).

\subsection{Tối thiểu hóa độ trải theo cách tăng dần}

Bộ giải được dựng một lần ở cận trên \(B\). Sau khi tìm mô hình có giá trị lớn
nhất \(\beta\), thêm
\[
  \neg g_{v,\beta}
  \qquad\forall v
\]
để buộc mọi nhãn không vượt quá \(\beta-1\). Vòng lặp nhảy theo nghiệm đương
nhiệm thật, không cần giảm từng đơn vị. Nếu lần gọi ở \(\beta-1\) trả \UNSAT,
mô hình tại \(\beta\) cùng cận dưới này chứng nhận tối ưu.

Đối xứng có thể khai thác theo tự đẳng cấu của đồ thị. Với chu trình \(C_n\),
sau khi dịch nhãn nhỏ nhất về \(1\), phép quay đưa đỉnh mang nhãn \(1\) về một
đỉnh neo cố định. Quy tắc neo như vậy phải dựa trên quỹ đạo của đồ thị.

\begin{resultbox}{Gán nhãn radio}
Trên 146 bài toán từ chín họ đồ thị, \textcite{thanh2026radio} báo 38 số radio
mới, khớp hoặc cải thiện \BKS trên 130 bài toán và chứng nhận 109 giá trị tối
ưu khi kết hợp các khung SAT với ILP. SAT thứ tự đặc biệt hiệu quả ở những họ
có đường kính tăng theo kích thước; ILP bổ sung tốt ở các họ có đường kính bị
chặn.
\end{resultbox}

\section[Gán nhãn an toàn k chính xác]{Gán nhãn an toàn \(k\) chính xác}

\subsection{Mô hình được sử dụng}

Trong mô hình của \textcite{thanh2026ksafe}, nhãn
\[
  L:V\rightarrow\set{1,\ldots,B}
\]
thỏa hai điều kiện:
\[
  L_u\ne L_v \qquad\forall u\ne v,
\]
\[
  \card{L_u-L_v}\ge k \qquad\forall uv\in E.
\]
Mục tiêu là tối thiểu \(\max_vL_v\). Điều kiện đầu là tính đơn ánh toàn cục:
mỗi nhãn được dùng nhiều nhất một lần.

\begin{designrule}{Quy ước của chương}
Chương giữ đúng mô hình đơn ánh ở trên. Vì vậy, khi \(k=1\), bài toán vẫn là
gán nhãn đơn ánh; đây không phải bài toán tô màu đồ thị chuẩn. Biến thể không đơn ánh là
một mô hình khác.
\end{designrule}

\subsection{Direct, forward order và reverse order encoding}

\paragraph{Direct Encoding (Direct Encoding, DE).}
\(K_{v,\ell}\) nói \(L_v=\ell\). Ràng buộc đúng một theo đỉnh và AMO theo nhãn
thực hiện phép gán đơn ánh. Mỗi cạnh cấm các cặp nhãn cách dưới \(k\).

\paragraph{Forward Order Encoding (Order Encoding xuôi, FE).}
\(X_{v,j}=1\iff L_v\ge j\). Giá trị chính xác xuất hiện ở điểm chuyển và kích
hoạt hai biên của khoảng \([j-k+1,j+k-1]\).

\paragraph{Reverse Order Encoding (Order Encoding ngược, RE).}
\(X_{v,j}=1\iff L_v\le j\). Ngữ nghĩa đối xứng với FE; hướng của chuỗi thứ tự
và các literal cận được đảo.

FE và RE có cùng tập phép gán nhãn nhưng có thể dẫn bộ giải đi theo các nhánh
khác nhau. Đối xứng phản xạ
\[
  L'_v=S+1-L_v
\]
giữ mọi khoảng cách và tính đơn ánh trên miền \([1,S]\). Có thể chọn một đỉnh
đại diện và đặt biến đại diện ở nửa dưới miền để giữ một đại diện của cặp phản xạ.

\begin{resultbox}{So sánh ba phép biểu diễn}
Trên 16 đồ thị Erdős--Rényi của \textcite{thanh2026ksafe}, FE và RE tìm nghiệm
cho toàn bộ tập và chứng nhận 13 giá trị tối ưu; DE tìm nghiệm cho 15 và chứng
nhận 8. Kết quả cho thấy biên thứ tự phù hợp với ràng buộc khoảng cách, còn
hướng FE/RE tạo khác biệt nhỏ hơn lựa chọn trực tiếp/thứ tự.
\end{resultbox}

\section{Phân bổ tần số với số giá trị nhỏ nhất}

\subsection{Miền rời rạc và hàm mục tiêu}

Mỗi máy phát \(v\) có miền tần số đã sắp
\[
  D_v=\set{D_{v,1}<D_{v,2}<\cdots<D_{v,m_v}}.
\]
Các cạnh mang khoảng cách nhiễu; một số máy phát có giá trị gán trước. Hàm mục
tiêu không phải tần số lớn nhất mà là số giá trị phân biệt được dùng:
\[
  \min\card{\set{f_v\mid v\in V}}.
\]
Do miền có lỗ và khác nhau giữa các máy phát, khoảng cấm phải được
tính theo giá trị thực của tần số rồi ánh xạ về chỉ số miền.

Trong hệ phân loại FAP, đây là mô hình số giá trị nhỏ nhất: hai nghiệm dùng cùng
số giá trị phân biệt có cùng giá trị mục tiêu dù độ trải khác nhau. Do đó chuỗi
thứ tự theo \emph{chỉ số miền cục bộ} chỉ phục vụ tính khả thi; biến \(A_a\)
theo \emph{giá trị toàn cục} mới phục vụ hàm mục tiêu. Trộn hai loại thứ tự này là một
lỗi mô hình phổ biến.

\begin{workedexample}{Cùng độ trải nhưng khác số giá trị}
Xét ba máy phát có miền chứa \(\{1,4,7\}\), và không có cạnh nhiễu giữa máy
phát thứ hai với thứ ba. Hai phương án
\[
  f=(1,4,7),\qquad f'=(1,7,7)
\]
đều có tần số lớn nhất \(7\), nhưng lần lượt dùng \(3\) và \(2\) giá trị phân
biệt. Tối thiểu hóa độ trải xem chúng ngang nhau; FAP số giá trị nhỏ nhất ưu
tiên \(f'\). Vì vậy biến hàm mục tiêu phải là \(A_1,A_4,A_7\), không phải chỉ
một ngưỡng ``tần số lớn nhất không vượt \(7\)''.
\end{workedexample}

\subsection{Tiền xử lý nhất quán cung}

Với cặp \((v,w)\), một giá trị \(a\in D_v\) phải có ít nhất một giá trị hỗ trợ
\(b\in D_w\) thỏa khoảng cách yêu cầu. Nếu không có, \(a\) bị loại. Quá trình
lặp đến \emph{fixpoint} (điểm bất động, tức không còn giá trị nào bị loại):
\[
  D_v\leftarrow
  \set{a\in D_v\mid
    \exists b\in D_w:\card{a-b}\ge d_{vw}}
\]
cho mọi cung ràng buộc. Tiền xử lý này không thay giá trị tối ưu vì chỉ loại giá
trị không thể xuất hiện trong bất kỳ nghiệm khả thi nào.

Đây là AC-3 trên các ràng buộc phân cách nhị phân. Khả năng suy diễn lan truyền của AC-3 bị giới hạn ở
hỗ trợ từng cạnh: một giá trị có thể có hỗ trợ riêng trên mọi đỉnh kề nhưng
không có hỗ trợ đồng thời toàn cục. Sau tiền xử lý, học mệnh đề SAT xử lý các
tổ hợp xung đột còn lại. Vì vậy nên báo kích thước miền tại điểm bất động và
thời gian AC riêng; không gán toàn bộ mức giảm thời gian chạy cho POSE.

\subsection{DSE và POSE}

Biểu diễn SAT trực tiếp dùng \(x_{v,q}=1\iff f_v=D_{v,q}\), ràng buộc đúng một
theo máy phát và xung đột từng cặp cho các cặp giá trị vi phạm.

Biểu diễn SAT thứ tự bộ phận thêm
\[
  g_{v,q}=1\iff
  \text{chỉ số tần số của }v\text{ không nhỏ hơn }q.
\]
Nếu \(x_{v,q}=1\), tần số \(D_{v,q}\) tạo khoảng cấm
\([D_{v,q}-d+1,D_{v,q}+d-1]\) trong miền \(D_w\). Hai biên
\(\lambda_w,\rho_w\) của miền rời rạc được đưa vào mệnh đề biên
\[
  \neg x_{v,q}
  \lor
  \neg g_{w,\lambda_w(q)}
  \lor
  g_{w,\rho_w(q)}.
\]
POSE vì thế loại sự phụ thuộc vào số giá trị nằm bên trong khoảng cấm.

\begin{figure}[tb]
\centering
\begin{tikzpicture}[satfig,x=13mm,y=9mm]
  \node[left,satlabel] at (0,1) {\(D_v\)};
  \foreach \i/\a in {1/1,2/4,3/7,4/10}{
    \ifnum\a=7
      \node[satbox,fill=PaleBlue,draw=Blue] at (\i,1) {\(\a\)};
    \else
      \node[satbox] at (\i,1) {\(\a\)};
    \fi
  }
  \node[left,satlabel] at (0,0) {\(D_w\)};
  \foreach \i/\a in {1/2,2/5,3/8,4/11}{
    \ifnum\i=1
      \node[satbox,fill=PaleTeal] at (\i,0) {\(\a\)};
    \else
      \ifnum\i=4
        \node[satbox,fill=PaleTeal] at (\i,0) {\(\a\)};
      \else
        \node[satbox,fill=PaleGold,draw=Gold] at (\i,0) {\(\a\)};
      \fi
    \fi
  }
  \draw[Gold,decorate,decoration={brace,amplitude=4pt}] (1.55,-.85)--(3.45,-.85)
    node[midway,below=4pt,satlabel,text=Gold]
      {bị cấm bởi \(f_v=7,d=4\)};
  \node[satlabel,align=left] at (7.4,.5)
    {\(\lambda_w=2\) tại giá trị \(5\)\\
     \(\rho_w=4\) tại giá trị \(11\)\\
     hỗ trợ còn lại: \(\{2,11\}\)};
\end{tikzpicture}
\caption{Ánh xạ khoảng giá trị \([4,10]\) vào hai biên của miền rời rạc
\(D_w=\{2,5,8,11\}\).}
\label{fig:fap-discrete-boundaries}
\end{figure}

\begin{workedexample}{Mệnh đề biên trên miền có lỗ}
Với các miền trong \cref{fig:fap-discrete-boundaries}, chọn giá trị thứ ba của
\(D_v\), tức \(x_{v,3}=1\) và \(f_v=7\). Khoảng cách \(d=4\) cấm
\([4,10]\) ở \(D_w\). Mệnh đề POSE là
\[
  \neg x_{v,3}\lor\neg g_{w,2}\lor g_{w,4}.
\]
Khi \(x_{v,3}=1\), \(\neg g_{w,2}\) cho phép chỉ số \(1\) (tần số \(2\));
\(g_{w,4}\) cho phép chỉ số \(4\) (tần số \(11\)). Hai giá trị \(5,8\) nằm
giữa bị loại mà không cần hai mệnh đề xung đột riêng.
\end{workedexample}

\subsection{Đếm số tần số phân biệt}

Với mỗi giá trị tần số toàn cục \(a\), tạo biến
\[
  A_a=1\iff
  \exists v:\ f_v=a.
\]
Liên kết \(x_{v,q}\rightarrow A_{D_{v,q}}\) bảo đảm mọi tần số được dùng đều
được đánh dấu. Hàm mục tiêu trở thành một ràng buộc đếm:
\[
  \min\sum_aA_a.
\]
Sequential Counter tăng dần được dựng một lần đến cận trên ban đầu. Sau khi mô hình
dùng \(q\) tần số, lần gọi kế tiếp kích hoạt cận \(q-1\) bằng các giả định. Lần
\UNSAT tại \(q-1\) cùng mô hình tại \(q\) chứng nhận tối ưu.

\begin{resultbox}{MO-FAP}
Trên 35 bài toán CALMA/DUTtest hoặc được mô hình hóa lại,
\textcite{dao2026mofap} báo cấu hình POSE kết hợp bộ đếm tăng dần chứng nhận
trạng thái khả thi hoặc vô nghiệm của toàn bộ tập thử. Nhất quán cung
giảm mạnh tổng thời gian của cấu hình này. Kết quả cho thấy vai trò bổ sung của
ba mô-đun: lọc miền, nén khoảng cấm và đếm giá trị phân biệt.
\end{resultbox}

\section{So sánh theo cấu trúc}

\begin{center}
\captionof{table}{So sánh cấu trúc Radio \(k\), \(k\)-Safe và MO-FAP.}
\label{tab:labeling-structure}
\begin{tabularx}{\textwidth}{@{}>{\bfseries\sffamily}p{.2\textwidth}YYY@{}}
\toprule
Đặc trưng & Radio \(k\) & \(k\)-Safe & MO-FAP\\
\midrule
Miền & \([1,B]\) liên tục & \([1,B]\) liên tục & \(D_v\) rời rạc\\
Khoảng cấm & Từ khoảng cách đồ thị & \(k\) trên cạnh & Từ nhiễu\\
Cặp bị ràng buộc & Mọi cặp có \(\Delta>0\) &
Cạnh và tính đơn ánh & Cạnh ràng buộc\\
Hàm mục tiêu & Nhãn lớn nhất nhỏ nhất & Nhãn lớn nhất nhỏ nhất & Số giá trị khác nhau nhỏ nhất\\
Lớp mục tiêu & Cận độ trải & Cận độ trải & Bộ đếm giá trị phân biệt\\
\bottomrule
\end{tabularx}
\end{center}

Ba bài toán chia sẻ mô-đun khả thi ở mức khoảng, nhưng lớp hàm mục tiêu khác
nhau. Radio và \(k\)-safe thay cận trên của nhãn. MO-FAP giữ các
giá trị số cố định và tối thiểu số giá trị được kích hoạt. Đây là lý do không
nên đồng nhất ``số giá trị nhỏ nhất'' với ``độ trải nhỏ nhất''.

\section{Đối chứng ngoài SAT và cận dưới}

Đối với độ trải gán nhãn, MIP dùng biến nhãn nguyên và hướng nhị phân cho trị
tuyệt đối; CP dùng ràng buộc khác nhau và khoảng cách; các phép dựng lý thuyết
đồ thị cho cận hoặc nghiệm trên từng họ đồ thị. Đối với FAP, MIP/sinh cột, tối
ưu ràng buộc, tìm kiếm cục bộ và \emph{tabu search} (tìm kiếm kiêng kỵ) là các đối
chứng đã được phát triển lâu
dài \parencite{aardal2007frequency}. Một đánh giá SAT
có giá trị nên phân biệt:
\begin{enumerate}
  \item bài toán mà SAT cải thiện nghiệm đương nhiệm;
  \item bài toán mà SAT đóng khoảng cách bằng chứng nhận \UNSAT;
  \item bài toán mà phương pháp khác cung cấp cận dưới/cận trên đầu vào;
  \item bài toán vô nghiệm do miền/giá trị gán trước, có thể phát hiện trước
        bước tối ưu.
\end{enumerate}

Với số giá trị nhỏ nhất, cận dưới cơ bản có thể đến từ cấu trúc đồ thị con đầy
đủ hoặc cấu trúc nhiễu
hoặc số giá trị bắt buộc bởi các giá trị gán trước; cận trên đến từ phương án
tần số kinh nghiệm. Bộ đếm giá trị phân biệt SAT tiếp nhận hai cận và lần lượt đóng
khoảng \([LB,UB]\).

\section{Bộ kiểm tra và chứng nhận}

Bộ kiểm tra khả thi đọc trực tiếp nghiệm:
\begin{itemize}
  \item radio: kiểm miền, đường đi ngắn nhất cho mọi cặp và mọi
        \(\Delta_{uv}>0\);
  \item \(k\)-safe: kiểm miền, tính đơn ánh và phân cách trên cạnh;
  \item MO-FAP: kiểm \(f_v\in D_v\), giá trị gán trước, mọi ràng buộc
        và số tần số phân biệt.
\end{itemize}
Bộ kiểm tra không đọc các biến thứ tự phụ. Điều này giúp bộ kiểm tra giữ kích thước nhỏ gọn và độc lập với
phép biểu diễn.

Chứng nhận tối ưu gồm một nghiệm làm chứng tại cận trên và bằng chứng cận dưới:
\UNSAT ở độ trải/số giá trị kề trước, hoặc một cận dưới toán học bằng cận trên.
Với các giả định tăng dần, có thể hiện thực hóa (reification) các giả định cuối thành mệnh đề đơn
vị để tạo một \CNF độc lập cho bộ kiểm tra chứng minh.

\begin{summarybox}
\begin{enumerate}
  \item Khoảng cấm là trừu tượng chung của cả ba bài toán.
  \item Trên miền liên tục, khoảng dùng hai ngưỡng trực tiếp; trên miền
        rời rạc, phải tìm biên theo giá trị.
  \item Tối thiểu hóa độ trải và tối thiểu hóa số giá trị phân biệt cần hai mô
        đun hàm mục tiêu khác nhau.
  \item Biến trực tiếp phục vụ lựa chọn chính xác và giải mã; biến thứ tự nén
        khoảng; Sequential Counter đếm số giá trị được dùng.
\end{enumerate}
\end{summarybox}

\subsection*{Câu hỏi nghiên cứu}
\begin{enumerate}
  \item Chứng minh \cref{lem:boundary} khi khoảng chạm một biên của miền.
  \item Với \(D_v=\set{10,30,40}\) và khoảng \([15,35]\), xác định hai
        literal biên của POSE.
\end{enumerate}
